\documentclass[11pt]{article}

\usepackage[utf8]{inputenc}
\usepackage[T1]{fontenc}
\usepackage[a4paper,margin=25mm]{geometry}
\usepackage{amsmath,amssymb}
\usepackage{graphicx}
\usepackage{subfigure}

% Images are stored in static/images, one level above this file.
\graphicspath{{../images/}}

\title{Latency Sensitivity Analysis in I-VIT}
\author{}
\date{}

\begin{document}
\maketitle

Latency sensitivity analysis is essential for identifying the operational boundaries of the system. Accordingly, we present a theoretical analysis below. Within the platform, the primary latency-sensitive link is the wireless communication between the physical vehicles and the cloud. 
Prior to the latency analysis, we first measured the step response of the scaled physical vehicle, as shown in Fig.~\ref{fig-physical-response-2}. 
The corresponding $2\%$ settling time is $15~\mathrm{ms}$.

\begin{figure}[h]
	%\vspace{-3mm}
	\centering
	\includegraphics[width=0.6\textwidth]{jieyue-sudu-partC.png}
	%\vspace{-2mm}
	\caption{Velocity step response of the scaled physical vehicle. 
    }
	\label{fig-physical-response-2}
	\vspace{-2mm}
\end{figure}


Considering a representative platooning scenario consistent with the experimental setup in the main text, we analyze an eight-vehicle CAV platoon consisting entirely of scaled physical vehicles.

I-VIT adopts a cloud-based control architecture in which the cloud aggregates the states of all vehicles, computes the complete CACC command for each following vehicle, and transmits the command to the corresponding scaled physical vehicle for execution. The measured latency of the motion capture system used to provide the physical-vehicle states to the cloud is $5~\mathrm{ms}$, which is sufficiently small to be neglected in the present analysis. The cloud is therefore assumed to have access to the real-time platoon states. Since the computation time required to generate the CACC-based control commands is also negligible, the following analysis focuses on a constant cloud-to-vehicle command delay $T_d$.
Consistent with the experimental setup in the main text, the platoon considered in this analysis adopts a predecessor–leader following (PLF) topology.

At time $t$, the cloud computes the command for the $i$-th following vehicle as
\[
u_i^{\mathrm{c}}(t) = k_s\left[s_i(t)-s^\ast\right] + k_l\left[v_1(t)-v_i(t)\right] + k_f\left[v_{i-1}(t)-v_i(t)\right],
\]
where $k_s=0.10$, $k_l=0.50$, $k_f=0.50$, and $s^\ast=7.89~\mathrm{m}$. Because the physical vehicle executes the cloud-generated command after the downlink delay, the actual control input applied at time $t$ is
\[
u_i(t) = u_i^{\mathrm{c}}(t-T_d).
\]
Equivalently,
\[
u_i(t) = k_s\left[s_i(t-T_d)-s^\ast\right] + k_l\left[v_1(t-T_d)-v_i(t-T_d)\right] + k_f\left[v_{i-1}(t-T_d)-v_i(t-T_d)\right].
\]
Accordingly, the delayed quantities in the above expression arise from the delayed execution of the complete cloud-generated command, rather than from delayed local sensing or delayed access to individual vehicle states.
The longitudinal actuation dynamics are approximated by a first-order system $\tau\dot{a}_i+a_i=u_i$. Given the measured 2\% settling time of $15~\mathrm{ms}$, the corresponding actuator time constant is $\tau \approx 0.003834~\mathrm{s}$.

% \vspace{1em}
\textbf{1) Closed-Loop Stability Analysis:} Closed-loop stability concerns whether the state of each vehicle eventually converges after the vehicle is subjected to a disturbance.
By linearizing around the equilibrium and applying the Laplace transform, the characteristic equation associated with each following vehicle is
\[
D(s,T_d) = s^2(\tau s+1) + e^{-sT_d} \left[ k_s+(k_l+k_f)s \right] = 0.
\]
To determine the admissible communication delay, we consider the first imaginary-axis crossing by setting $s=j\omega$. The magnitude condition of the characteristic equation is
\[
\omega^4 \left( 1+\tau^2\omega^2 \right) = k_s^2+(k_l+k_f)^2\omega^2.
\]
Substituting the adopted parameters and solving this equation yields the cross-over frequency $\omega_{\mathrm{cl}} \approx 1.005~\mathrm{rad/s}$. The corresponding closed-loop delay margin is
\[
T_{d,\mathrm{cl}} = \frac{ \tan^{-1} \left( \frac{(k_l+k_f)\omega_{\mathrm{cl}}}{k_s} \right) - \tan^{-1} \left( \tau\omega_{\mathrm{cl}} \right) }{ \omega_{\mathrm{cl}} } \approx 1.46~\mathrm{s}.
\]
Therefore, the closed-loop system of any individual vehicle remains internally stable when
\[
\boxed{ T_d < 1.46~\mathrm{s} }.
\]
Fig.~\ref{fig-cacc-closed-loop} shows the Bode phase plots of the open-loop transfer function, which are used to evaluate closed-loop stability under different cloud-to-vehicle command delays. The delay introduces additional phase lag without changing the gain crossover frequency, and the phase reaches $-180^\circ$ at $T_d\approx1.46~\mathrm{s}$, defining the closed-loop stability limit.
\begin{figure*}[h!]
	%\vspace{1mm}
	\centering
	\subfigure[Closed-loop stability]
	{\includegraphics[width=0.45\textwidth]{CACC_closed_loop.png}
		\label{fig-cacc-closed-loop}}
	\hspace{4mm}
	\subfigure[String stability]
	{\includegraphics[width=0.45\textwidth]{CACC_string_stability.png}
		\label{fig-CACC-string-stability}}
	\vspace{-1mm}
	\caption{Frequency-domain analysis of (a) closed-loop stability and (b) spacing-error string stability of the CACC platoon under different cloud-to-vehicle communication delays.}
	\label{fig-CACC-stability}
\end{figure*}

\textbf{2) String-Stability Analysis:}
String stability requires that spacing disturbances are attenuated rather than amplified along the platoon. The spacing-error propagation between consecutive vehicles is governed by $E_{i+1}(s) = A(s,T_d)E_i(s)$, where the transfer function is
\[
A(s,T_d) = \frac{ e^{-sT_d}(k_s+k_fs) }{ s^2(\tau s+1) + e^{-sT_d} \left[ k_s+(k_l+k_f)s \right] }.
\]
Accordingly, spacing-error string stability requires $\sup_{\omega\geq0} \left| A(j\omega,T_d) \right| \leq 1$. Because $A(0,T_d)=1$, the critical delay corresponds to the first nonzero frequency at which the magnitude reaches unity. This critical boundary is obtained by jointly solving
\[
\left| A(j\omega,T_d) \right| = 1 \quad \text{and} \quad \frac{ \partial \left| A(j\omega,T_d) \right| }{ \partial\omega } = 0.
\]
Solving these equations yields the critical frequency $\omega_{\mathrm{ss}} \approx 1.306~\mathrm{rad/s}$ and the corresponding string-stability delay margin
\[
T_{d,\mathrm{ss}} \approx 0.87~\mathrm{s}.
\]
When $0.87~\mathrm{s} < T_d < 1.46~\mathrm{s}$, individual vehicles remain stable, but disturbances may be amplified towards the tail of the platoon, since $E_8(s)/E_2(s) = A^6(s,T_d)$. 
Therefore, spacing-error string stability for the entire platoon requires
\[
\boxed{ T_d \leq 0.87~\mathrm{s} }.
\]
Fig.~\ref{fig-CACC-string-stability} shows the Bode magnitude plots of the inter-vehicle spacing-error propagation transfer function. The propagation gain reaches the $0~\mathrm{dB}$ boundary at $T_d\approx0.87~\mathrm{s}$, beyond which spacing errors are amplified along the platoon. Therefore, spacing-error string stability imposes a more restrictive delay requirement than closed-loop stability.

\begin{figure*}[h!]
	%\vspace{1mm}
	\centering
	\subfigure[Velocity profiles]
	{\includegraphics[width=0.45\textwidth]{delay-1500ms-speed-partC.png}
		\label{fig-exp-delay-velocity}}
	\hspace{4mm}
	\subfigure[Inter-vehicle distance profiles]
	{\includegraphics[width=0.45\textwidth]{delay-1500ms-distance-partC.png}
		\label{fig-exp-delay-distance}}
	\vspace{-1mm}
	\caption{Velocity and inter-vehicle spacing profiles of all virtual vehicles in the communication-delay sensitivity experiment. The vehicles are controlled by delayed CACC commands issued by the cloud, with a delay of $1.5~\mathrm{s}$, which exceeds the theoretically derived closed-loop stability limit of $1.46~\mathrm{s}$ for the platoon.}
	\label{fig-exp-delay-result}
\end{figure*}

The measured communication latency between the scaled physical vehicles and the cloud on the I-VIT platform follows a Gaussian distribution with a mean of $1.33~\mathrm{ms}$ and a standard deviation of $0.66~\mathrm{ms}$. 
This latency is substantially lower than the derived thresholds and therefore fully satisfies the requirements for conducting CACC platooning experiments in terms of both closed-loop platoon stability and spacing-error string stability.

However, when the communication latency exceeds a certain threshold, the platform can no longer operate stably. 
Taking the classical CACC platooning test as an example, 
we conducted platoon simulations under different communication delays using a longitudinal dynamics model calibrated to the response characteristics of the scaled physical vehicles. As shown in Fig.~\ref{fig-exp-delay-result}, when the communication latency between the cloud and the vehicles reaches $1.50~\mathrm{s}$, exceeding the closed-loop stability threshold of $1.46~\mathrm{s}$, a velocity disturbance introduced by the head vehicle causes both the vehicle velocities and inter-vehicle spacings to diverge throughout the platoon. Consequently, the platform can no longer operate stably.

\end{document}
